\section{Введение}

Учебная ознакомительная практика (компьютерная практика) проводится
с целью закрепления навыков программирования на языке высокого
уровня~C, изучения методов решения математических и геометрических
задач, а также приобретения опыта документирования программных
решений.

\subsection{Цель практики}

Закрепить навыки программирования на языке~C путём самостоятельной
разработки, тестирования и документирования 15 программ, охватывающих
основные разделы курса: линейные, разветвляющиеся и циклические
алгоритмы, операции над массивами и матрицами, линейный поиск,
арифметику, геометрию и теорию множеств, линейную алгебру,
обработку символьной информации, аналитическую геометрию,
кривые и поверхности второго порядка.

\subsection{Задачи практики}

\begin{enumerate}
  \item Изучить постановки 15 задач (по одной из каждой темы методических
        указаний) согласно индивидуальному варианту~4.
  \item Для каждой задачи привести математическое обоснование решения
        и словесное описание алгоритма.
  \item Разработать блок-схему алгоритма средствами \LaTeX{} (TikZ).
  \item Реализовать решение на языке~C с пояснительными комментариями.
  \item Подготовить и протестировать программу на контрольных данных.
  \item Оформить отчёт по практике в соответствии с требованиями
        к структуре и форматированию документа.
\end{enumerate}

\subsection{Перечень выполненных задач (вариант~4)}

\begin{center}
\begin{longtable}{cl}
\toprule
№ & Тема (задача~4 варианта) \\
\midrule
1  & Линейные алгоритмы --- продолжительность временного промежутка \\
2  & Разветвляющиеся алгоритмы --- тип четырёхугольника \\
3  & Циклические алгоритмы --- площадь пересечения прямоугольников \\
4  & Операции над массивами --- среднее значение и дисперсия \\
5  & Векторы и матрицы --- проекция трёхмерного куба \\
6  & Линейный поиск --- верхняя огибающая линейных функций \\
7  & Арифметика --- округление до 3 значащих цифр \\
8  & Геометрия и теория множеств --- подмножества с суммой в диапазоне \\
9  & Линейная алгебра --- вычисление минора матрицы \\
10 & Обработка символьной информации --- самое длинное слово и фраза \\
11 & Аналитическая геометрия --- угол между прямыми \\
12 & Кривые второго порядка --- вершины повёрнутого прямоугольника \\
13 & Графическое решение систем уравнений --- метод Крамера \\
14 & Плоскость в пространстве --- расстояние и проекция точки \\
15 & Поверхность второго порядка --- положение точки относительно эллипсоида \\
\bottomrule
\end{longtable}
\end{center}

\subsection{Среда разработки}

Все программы написаны на языке~C (стандарт C11) и скомпилированы
компилятором \texttt{gcc} с флагами \texttt{-Wall~-Wextra} (без
предупреждений).  Блок-схемы построены средствами пакета \texttt{TikZ}
системы \LaTeX{} по образцам, приведённым в методических материалах
практики.
